\documentclass[a4paper,11pt]{ltjsarticle}
\usepackage{base}
\title{}
\author{}
\date{}
\newcommand{\printheader}[2]{
\begin{tikzpicture}[remember picture, overlay]
\node[yshift=-2.5cm, anchor=north] at (current page.north) {
\begin{tikzpicture}
\fill[gray!20] (0,0) rectangle (\textwidth, 2cm);
\fill[gray!80] (0,0) rectangle (0.2cm, 2cm);
\draw[gray!80, thick] (0,0) -- (	\textwidth, 0);
\node[anchor=west, text width=\textwidth-1cm, inner xsep=1cm] at (0, 1.25cm) {
\parbox[b]{\linewidth}{
{\color{gray!50!black}\bfseries #1} \par
\vspace{0.2em}
{\huge\bfseries #2}
}
};
\end{tikzpicture}
};
\end{tikzpicture}
\vspace{0.5cm}
}
\begin{document}
\printheader{11月30日}{積分法・数列と極限① 解答}

\begin{toi}
(1) 適当に倍角の公式で変形する．
\hfill \textbf{答} \quad $\boldsymbol{\sin x=\dfrac{2t}{1+t^2}, \quad \cos x=\dfrac{1-t^2}{1+t^2}, \quad \tan x=\dfrac{2t}{1-t^2}}$

\vspace{1em}

(2) $t=\tan\dfrac{x}{2}$とおくと， $dx=\dfrac{2}{1+t^2}dt$．
\[
\begin{aligned}
\int \frac{1}{3\sin x +4\cos x}dx &= \int \frac{1}{3\left(\frac{2t}{1+t^2}\right) + 4\left(\frac{1-t^2}{1+t^2}\right)} \cdot \frac{2}{1+t^2}dt \\
&= \int \frac{2}{6t + 4(1-t^2)} dt = \int \frac{-1}{2t^2-3t-2} dt \\
&= \int \frac{-1}{(2t+1)(t-2)} dt = \frac{1}{5}\int \left(\frac{2}{2t+1} - \frac{1}{t-2}\right) dt \\
&= \frac{1}{5}(\log|2t+1| - \log|t-2|) + C \\
&= \frac{1}{5}\log\left|\frac{2t+1}{t-2}\right| + C\\
&=\dfrac{1}{5}\log\left|\dfrac{2\tan\frac{x}{2}+1}{\tan\frac{x}{2}-2}\right|+C
\end{aligned}
\]
\hfill \textbf{答} \quad $\boldsymbol{\dfrac{1}{5}\log\left|\dfrac{2\tan\frac{x}{2}+1}{\tan\frac{x}{2}-2}\right|+C}$
\end{toi}

\begin{toi}
(1) $t=\sqrt{x}$とおくと，$x=t^2, dx=2tdt$．$x:1\to9$のとき$t:1\to3$．
\[
\begin{aligned}
\int_1^9 \sqrt{x}e^{\sqrt{x}}dx &= \int_1^3 t e^t \cdot 2t dt = 2\int_1^3 t^2 e^t dt \\
&= 2\left( [t^2 e^t]_1^3 - \int_1^3 2t e^t dt \right) \\
&= 2(9e^3 - e) - 4( [te^t]_1^3 - \int_1^3 e^t dt ) \\
&= 18e^3 - 2e - 4(3e^3 - e - (e^3 - e)) = 10e^3 - 2e
\end{aligned}
\]
\hfill \textbf{答} \quad $\boldsymbol{2e(5e^2-1)}$

\vspace{0.5em}

(2) $\dfrac{1}{12}$公式を使ってもいいが，覚えるのがめんどうなので部分積分でいいと思う．
\[
\int_\alpha^\beta (x-\alpha)(x-\beta)^2dx = \left[\frac{(x-\alpha)(x-\beta)^3}3\right]_\alpha^\beta-\frac13\int_\alpha^\beta(x-\beta)^4dx=\frac{(\alpha-\beta)^4}{12} \]
 　\hfill\textbf{答} \quad $\boldsymbol{\dfrac{(\alpha-\beta)^4}{12}}$

\vspace{0.5em}

(3) 部分積分法を用いる．
\[
\begin{aligned}
\int_1^e x(\log x)^2dx &= \left[\frac{x^2}{2}(\log x)^2\right]_1^e - \int_1^e \frac{x^2}{2} \cdot 2\log x \cdot \frac{1}{x} dx \\
&= \frac{e^2}{2} - \int_1^e x\log x dx = \frac{e^2}{2} - \left( \left[\frac{x^2}{2}\log x\right]_1^e - \int_1^e \frac{x}{2}dx \right) \\
&= \frac{e^2}{2} - \frac{e^2}{2} + \left[\frac{x^2}{4}\right]_1^e = \frac{e^2-1}{4}
\end{aligned}
\]
\hfill \textbf{答} \quad $\boldsymbol{\dfrac{e^2-1}{4}}$

\vspace{0.5em}

(4) 
\[
\begin{aligned}
\int_0^\pi x^2\sin x dx &= [-x^2\cos x]_0^\pi - \int_0^\pi 2x(-\cos x)dx \\
&= \pi^2 + 2\int_0^\pi x\cos x dx \\
&= \pi^2 + 2\left( [x\sin x]_0^\pi - \int_0^\pi \sin x dx \right) \\
&= \pi^2 + 2(0 - [-\cos x]_0^\pi) = \pi^2 - 4
\end{aligned}
\]
\hfill \textbf{答} \quad $\boldsymbol{\pi^2-4}$
\end{toi}

\begin{toi}
(1) 漸化式を立てる．
$n+1$回目に奇数回となるのは，「$n$回目奇数かつ$n+1$回目起こらない」または「$n$回目偶数かつ$n+1$回目起こる」場合である．
\[
a_{n+1} = a_n(1-p) + (1-a_n)p = (1-2p)a_n + p
\]
特性方程式 $\alpha = (1-2p)\alpha + p$ を解いて $\alpha = \frac{1}{2}$．
数列 $\left\{a_n - \frac{1}{2}\right\}$ は公比 $1-2p$ の等比数列．初項は $a_1=p$ より $a_1-\frac{1}{2} = p-\frac{1}{2}$．
\[
a_n - \frac{1}{2} = \left(p-\frac{1}{2}\right)(1-2p)^{n-1} = -\frac{1}{2}(1-2p)^n
\]
\hfill \textbf{答} \quad $\boldsymbol{a_n = \dfrac{1}{2}\{1-(1-2p)^n\}}$

\vspace{1em}

(2) $0<p<1$ より $-1 < 1-2p < 1$ なので $\lim_{n\to\infty}(1-2p)^n = 0$．
\hfill \textbf{答} \quad $\boldsymbol{\dfrac{1}{2}}$
\end{toi}
\newpage
\begin{toi}
(1) $I_1 = \int_1^e \log x dx = [x\log x - x]_1^e = (e-e) - (0-1) = 1$
\hfill \textbf{答} \quad $\boldsymbol{1}$

\vspace{0.5em}

(2) 
\[
\begin{aligned}
I_{n+1} &= \int_1^e 1 \cdot (\log x)^{n+1} dx \\
&= [x(\log x)^{n+1}]_1^e - \int_1^e x \cdot (n+1)(\log x)^n \cdot \frac{1}{x} dx \\
&= e - (n+1)I_n
\end{aligned}
\]
\hfill \textbf{答} \quad $\boldsymbol{I_{n+1} = e - (n+1)I_n}$

\vspace{0.5em}

(3) 
$I_1 = 1$ ,~~~$I_2 = e - 2(1) = e-2$ ,~~~$I_3 = e - 3(e-2) = -2e+6$ ,~~~$I_4 = e - 4(-2e+6) = 9e-24$\\
　\hfill \textbf{答} \quad $\boldsymbol{9e-24}$
\end{toi}

\begin{toi}
(1) 普通に計算するだけ．
\hfill \textbf{答} \quad $\boldsymbol{\dfrac{1}{2}\left(\dfrac{1}{2k-1} - \dfrac{1}{2k+1}\right)}$

\vspace{1em}

(2) 
\[
\begin{aligned}
S_n &= \frac{n}{2} \sum_{k=n}^{2n} \left(\frac{1}{2k-1} - \frac{1}{2k+1}\right) \\
&= \frac{n}{2} \left\{ \left(\frac{1}{2n-1} - \frac{1}{2n+1}\right) + \cdots + \left(\frac{1}{4n-1} - \frac{1}{4n+1}\right) \right\} \\
&= \frac{n}{2} \left( \frac{1}{2n-1} - \frac{1}{4n+1} \right) = \frac{1}{2} \left( \frac{1}{2-1/n} - \frac{1}{4+1/n} \right)
\end{aligned}
\]
$n\to\infty$のとき $S_n\to\dfrac{1}{2}\left(\dfrac{1}{2} - \dfrac{1}{4}\right) = \dfrac{1}{8}$
\hfill \textbf{答} \quad $\boldsymbol{\dfrac{1}{8}}$
\end{toi}

\begin{toi}
ガウス記号の性質 $x-1 < [x] \leqq x$ を用いてはさみうちの原理を適用する．

\noindent
\begin{minipage}{0.33\linewidth}
(1) $\dfrac{n}{3}-1 < \left[\dfrac{n}{3}\right] \leqq \dfrac{n}{3}$ より\\
$\dfrac{1}{3}-\dfrac{1}{n} < \dfrac{[n/3]}{n} \leqq \dfrac{1}{3}$\\
極限は $\boldsymbol{1/3}$
\end{minipage}
\begin{minipage}{0.33\linewidth}
(2) $\sqrt{n}-1 < [\sqrt{n}] \leqq \sqrt{n}$ より\\
$1-\dfrac{1}{\sqrt{n}} < \dfrac{[\sqrt{n}]}{\sqrt{n}} \leqq 1$\\
極限は $\boldsymbol{1}$
\end{minipage}
\begin{minipage}{0.33\linewidth}
(3) $10^n\pi-1 < [10^n\pi] \leqq 10^n\pi$ より\\
$\pi-\dfrac{1}{10^n} < \dfrac{[10^n\pi]}{10^n} \leqq \pi$\\
極限は $\boldsymbol{\pi}$
\end{minipage}
\end{toi}

\begin{toi}
(1)二項定理より $n \geqq 2$ のとき
\[
(1+h)^n = 1 + {}_n\mathrm{C}_1 h + {}_n\mathrm{C}_2 h^2 + \cdots + h^n
\]
$h>0$ より第3項以降は正であるから
\[
(1+h)^n \geqq 1 + nh + \frac{n(n-1)}{2}h^2
\]
\hfill \textbf{(証明終)}

\vspace{1em}

(2) (1)より $a^n = (1+h)^n > \dfrac{n(n-1)}{2}h^2$．
逆数をとり $n$ を掛けると
\[
0 < \frac{n}{a^n} < \frac{n}{\frac{n(n-1)}{2}h^2} = \frac{2}{(n-1)h^2}
\]
$n\to\infty$ で右辺は0に収束するため，はさみうちの原理より
\hfill \textbf{答} \quad $\boldsymbol{\lim_{n\to\infty}\dfrac{n}{a^n}=0}$
\end{toi}

\end{document}